% ============================================================
% 279_T0_Casimir_CMB_H0_En_ch.tex
% Step-by-step: Casimir <-> CMB <-> H0 via the ξ scale
% Sources: Doc. 091, 165, 025/061, 166, 026; status: Doc. 190
% ============================================================

\section*{What this is about}

Three phenomena that at first have nothing to do with one another: the
\emph{Casimir effect} (vacuum force between plates, $\sim\mu$m), the
\emph{cosmic microwave background} (CMB, $T_\text{CMB}=2{,}725$\,K) and the
\emph{Hubble constant} $H_0$ (cosmic expansion rate). In FFGFT all three hang
on \emph{one} vacuum length scale $L_\xi\approx100\,\mu$m and on the same
exponent $41/4$ of the geometry parameter $\xi=4/30000$. This note lays the
connection open step by step (brought together in Doc.~166 from Doc.~091, 165,
025/061; status assessment per Doc.~190).

\medskip
Throughout: $\xi = \dfrac{4}{30000} = 1{,}3333\times10^{-4}$.

% ============================================================
\section{Preliminary remark: why \enquote{Hubble constant} is misleading}
% ============================================================

Before we use $H_0$, a terminological clarification -- because the usual
designation \emph{Hubble constant} is, from the FFGFT view, misleading on two
counts.

\textbf{First: not a constant.} Already in standard cosmology $H_0$ is not
constant but the \emph{present-day value} of the time-dependent Hubble
parameter $H(t)=\dot a/a$. \enquote{Constant} there means only spatially
homogeneous, not fixed in time. The name is thus already an established
imprecision there.

\textbf{Second (and sharper): in FFGFT there is no expansion rate at all.}
The static $\xi$ universe has no scale factor $a(t)$ and no time-varying
$H(t)$; the redshift is fractal path lengthening, not an expansion effect
(Doc.~190, R14a/R15). There is therefore no \enquote{rate} whose present-day
value one could take as $H_0$. The expansion reading is a $\Lambda$CDM pipeline
interpretation (R18), not an FFGFT statement.

\textbf{What remains physically is a length.} The robust, model-neutral
quantity is the \emph{Hubble length}
\begin{equation}
  R_H \;=\; \frac{c}{H_0},
  \label{eq:RH}
\end{equation}
and $H_0 = c/R_H$ is merely this length expressed, via the SI conversion factor
$c$, as an inverse time. $R_H$ is a \emph{reference length} -- the upper scale
of the cosmological torus sector -- entirely analogous to the Planck length
$\ell_P$ or to the system-specific bounds $R = \hbar/mc$. It is \emph{not} the
size of the universe and \emph{not} the observable horizon (the latter
$\approx 3\,R_H$ in $\Lambda$CDM); only a lower bound is measurable (Doc.~190,
R36).

\begin{foundation}[Terminology convention]
In this note the cosmic scale is conceived primarily as a \emph{reference
length} $R_H$, not as a \enquote{constant} $H_0$. Where $H_0$ appears it is to
be read as $H_0 = c/R_H$ -- an SI restatement of the reference length, not a
fundamental constant of nature. This is consistent with $H_0$ as an emergent,
decompactified quantity (R16) and with the status of $R_H$ as a shared
empirical reference point (R39). Precisely for this reason the following
derivation works throughout with \emph{lengths} ($L_\xi$, $R_H$) and carries
$H_0$ only as their conversion.
\end{foundation}

% ============================================================
\section{Starting point: three relations from three documents}
% ============================================================

\begin{foundation}[The three building blocks]
\begin{enumerate}
  \item \textbf{Casimir/vacuum (Doc.~091):} the CMB energy density as a T0
        vacuum density,
        \begin{equation}
          \rho_\text{CMB} = \frac{\xi\,\hbar c}{L_\xi^4}.
          \label{eq:rhocmb}
        \end{equation}
  \item \textbf{$H_0$ ratio (Doc.~165):}
        \begin{equation}
          \frac{\hbar H_0}{m_e c^2} = \frac{\pi}{2}\,\xi^{10}.
          \label{eq:h0}
        \end{equation}
  \item \textbf{CMB temperature (Doc.~025/061):} $T_\text{CMB}=2{,}72548$\,K,
        brought in via the Stefan-Boltzmann law (Step~1).
\end{enumerate}
\end{foundation}

\noindent Relations~\eqref{eq:rhocmb} and~\eqref{eq:h0} were developed
independently. The question: are they two aspects of the same $\xi$ structure?

% ============================================================
\section{Step 1: Casimir $\leftrightarrow$ CMB via the scale $L_\xi$}
% ============================================================

\subsection*{1a -- Equating the two energy densities}

The standard Casimir energy density between two plates at separation $d$ is
\begin{equation}
  \rho_\text{Casimir}(d) = \frac{\pi^2\,\hbar c}{240\,d^4}.
  \label{eq:rhocas}
\end{equation}
At exactly the separation $d=L_\xi$ it relates to the vacuum density
\eqref{eq:rhocmb} as
\begin{equation}
  \boxed{\;\frac{\rho_\text{Casimir}(L_\xi)}{\rho_\text{CMB}}
  = \frac{\pi^2\hbar c/(240\,L_\xi^4)}{\xi\hbar c/L_\xi^4}
  = \frac{\pi^2}{240\,\xi}\approx 308\;}
  \label{eq:ratio}
\end{equation}
Remarkable: inserting \eqref{eq:rhocmb} into \eqref{eq:ratio}, $\xi$ cancels and
one recovers \emph{exactly} the standard Casimir formula. The T0 notation is
thus consistent with the textbook Casimir formula; $L_\xi$ is the scale at which
the vacuum and Casimir densities are of the same order.

\subsection*{1b -- $L_\xi$ from the CMB temperature}

The CMB density is thermal (Stefan-Boltzmann):
\begin{equation}
  \rho_\text{CMB} = \frac{\pi^2}{15}\,\frac{(k_B T_\text{CMB})^4}{(\hbar c)^3}.
  \label{eq:sb}
\end{equation}
Equating \eqref{eq:rhocmb} and \eqref{eq:sb} and solving for $L_\xi$:
\begin{equation}
  \frac{\xi\hbar c}{L_\xi^4} = \frac{\pi^2}{15}\frac{(k_BT_\text{CMB})^4}{(\hbar c)^3}
  \;\Longrightarrow\;
  \boxed{\;L_\xi = \left[\frac{15\,\xi}{\pi^2}\right]^{1/4}\frac{\hbar c}{k_B T_\text{CMB}}\;}
  \label{eq:Lxi}
\end{equation}

\begin{keyresult}[Casimir exponent $1/4$]
Numerically: $L_\xi = 100{,}24\,\mu$m (0.11\,\% from the scale determined
directly from $\rho_\text{CMB}$). Decisive for what follows: from
\eqref{eq:Lxi}, $\xi$ enters $L_\xi$ \emph{with the power $1/4$},
\[
  L_\xi \;\propto\; \xi^{1/4}.
\]
\end{keyresult}

% ============================================================
\section{Step 2: $H_0$ as a $\xi$ power}
% ============================================================

Relation~\eqref{eq:h0} solved for $H_0/c$ (with the reduced Compton wavelength
$\bar\lambda_e=\hbar/m_ec$):
\begin{equation}
  \frac{H_0}{c} = \frac{\pi}{2}\,\xi^{10}\,\frac{m_e c^2}{\hbar c}
  = \frac{\pi}{2}\,\frac{\xi^{10}}{\bar\lambda_e}.
  \label{eq:h0overc}
\end{equation}

\begin{keyresult}[$H_0$ exponent $10$]
Inserted, this gives $H_0 \approx 66{,}8$\,km/s/Mpc -- within $0{,}9\,\%$ of the
Planck value $67{,}4$\,km/s/Mpc, without free parameters. Here $\xi$ enters with
the power $10$:
\[
  H_0 \;\propto\; \xi^{10}.
\]
\end{keyresult}

% ============================================================
\section{Step 3: Multiplication -- the bridge formula}
% ============================================================

Now the product of the two scales $L_\xi$ \eqref{eq:Lxi} and $H_0/c$
\eqref{eq:h0overc}. The powers of $\xi$ add, $\hbar c$ cancels:
\begin{align}
  L_\xi\cdot\frac{H_0}{c}
  &= \left[\frac{15\xi}{\pi^2}\right]^{1/4}\frac{\hbar c}{k_BT_\text{CMB}}
     \;\cdot\; \frac{\pi}{2}\,\frac{m_ec^2\,\xi^{10}}{\hbar c}
     \nonumber\\[4pt]
  &= \frac{\pi}{2}\left[\frac{15}{\pi^2}\right]^{1/4}
     \frac{m_ec^2}{k_BT_\text{CMB}}\;\cdot\;\xi^{1/4}\,\xi^{10}
     \nonumber\\[4pt]
  &= \frac{\pi}{2}\left[\frac{15}{\pi^2}\right]^{1/4}
     \frac{m_ec^2}{k_BT_\text{CMB}}\;\cdot\;\xi^{41/4}.
  \label{eq:bridgederiv}
\end{align}

\begin{keyresult}[Connection formula Casimir $\leftrightarrow$ $H_0$]
\begin{equation}
  L_\xi\cdot\frac{H_0}{c}
  = \underbrace{\frac{\pi}{2}\left(\frac{15}{\pi^2}\right)^{1/4}}_{K_\text{geo}=1{,}7441}
  \cdot
  \underbrace{\frac{m_ec^2}{k_BT_\text{CMB}}}_{=\,2{,}176\times10^{9}}
  \cdot\;\xi^{41/4}.
  \label{eq:bridge}
\end{equation}
Numerical check: $L_\xi\cdot H_0/c = 7{,}241\times10^{-31}$ -- computed directly
and via \eqref{eq:bridge} algebraically identical (to $10^{-14}\,\%$). The
derivation is \emph{exact} -- no approximation, no free parameters
\emph{between} the building blocks.
\end{keyresult}

% ============================================================
\section{The exponent $41/4$ and its decomposition}
% ============================================================

The exponent is the sum of the two individual exponents from Steps 1 and 2:
\begin{equation}
  \boxed{\;\frac{41}{4}
  = \underbrace{10}_{\substack{H_0\text{ exponent}\\ \eqref{eq:h0}}}
  + \underbrace{\frac14}_{\substack{\text{Casimir exponent}\\ \eqref{eq:Lxi}}}\;}
  \label{eq:decomp}
\end{equation}
The same exponent $41/4$ appears independently in Doc.~026 as
$E_H = E_0\,\xi^{41/4}$ (Hubble energy scale). The decomposition $41/4=10+1/4$ is
therefore not constructed by chance: the $10$ carries the coupling electron
mass~$\to$~Hubble scale, the $1/4$ carries the Casimir scaling of the vacuum
energy.

% ============================================================
\section{Where does the SI conversion factor sit?}
% ============================================================

A legitimate counter-question: we are operating here at the
\emph{decompactified} level of the formulation -- $H_0$ exists only after
decompactification (Doc.~190, R16: $H_0 = E_H/\hbar$ contains the $\hbar$ that
becomes real only then), and $L_\xi$ is a \emph{physical} length ($\approx
100\,\mu$m) built from $\hbar$, $c$, $k_B$. At this level an SI conversion
factor \emph{must} appear. It does -- it merely sits hidden in two places, and
in the exponent notation in a third.

\subsection*{Place 1: the factor cancels (legitimate)}

Both scales are built with the same $\hbar c$:
\begin{equation}
  L_\xi \;\propto\; \frac{\hbar c}{k_B T_\text{CMB}},
  \qquad
  \frac{H_0}{c} \;=\; \frac{\pi}{2}\,\xi^{10}\,\frac{m_e c^2}{\hbar c}.
\end{equation}
In the product $\hbar c$ drops out -- because $L_\xi\cdot H_0/c$ is at its core
a \emph{length ratio} ($L_\xi/\bar\lambda_e$ with the reduced Compton wavelength
$\bar\lambda_e=\hbar/m_ec$), and the ratio of two lengths is unit-free. This is
no trick: the same SI conversion $\hbar c$ builds \emph{both} lengths, so it
must cancel in the ratio.

\subsection*{Place 2: the factor survives (the actual anchor)}

What remains after the cancellation is \emph{not} purely geometric:
\begin{equation}
  \boxed{\;\frac{m_e c^2}{k_B T_\text{CMB}} \;=\; 2{,}176\times10^{9}\;}
\end{equation}
This contains $c^2$ (mass $\to$ energy) and $k_B$ (Kelvin $\to$ energy) -- and
it is a \emph{ratio of two measured quantities}: electron rest energy to CMB
temperature. \emph{That} is the SI / empirical conversion factor. The bridge
formula \eqref{eq:bridge} is thereby geometric \emph{form}
($K_\text{geo}\,\xi^{41/4}$) \textbf{times} an SI anchor
($m_e c^2/k_B T_\text{CMB}$) -- \emph{not} a pure geometric identity.

\subsection*{Place 3: in the exponent -- $41/4$ vs.\ $63/4$}

The factor shows up most cleanly in the exponent itself. The $41/4$ is
referenced to the electron / CMB energy scale $E_0$, \emph{not} to the Planck
scale. Planck-referenced (Doc.~190, R35) the exponent is $63/4$, and
\begin{equation}
  \boxed{\;\frac{63}{4} \;=\; \frac{41}{4} \;+\; \frac{11}{2}\;}
  \label{eq:exp-split}
\end{equation}
The $11/2$ \emph{is} exactly the SI scale conversion factor
$\ell_P \leftrightarrow E_0$ (Planck length vs.\ electron / CMB scale). Put
differently: if one writes $\xi^{41/4}$ and calls it \enquote{purely
geometric}, one has tacitly absorbed $\ell_P$ into the exponent -- exactly the
guardrail R35 (\enquote{the SI conversion factor $\ell_P$ must not be absorbed
into the exponent}). Which reference scale one makes explicit is a \emph{choice};
the factor never disappears, it only changes location:
\begin{center}
{\small
\renewcommand{\arraystretch}{1.3}
\begin{tabular}{p{4.4cm}p{4.4cm}p{4.4cm}}
\toprule
\textbf{Reference scale} & \textbf{Exponent} & \textbf{SI factor explicit as} \\
\midrule
$E_0$ (electron/CMB) & $41/4$ & ratio $m_ec^2/k_BT_\text{CMB}$ \\
$\ell_P$ (Planck)    & $63/4 = 41/4 + 11/2$ & extra exponent $11/2$ \\
\bottomrule
\end{tabular}
}
\end{center}

\subsection*{Why we are necessarily decompactified here}

This is not a matter of style: $H_0$ \emph{does not exist} in the compact $T^4$
form -- it contains the $\hbar$ that becomes real only upon decompactification
(R16) -- and $L_\xi$ is a concrete $\mu$m length that cannot be formulated
without $\hbar, c, k_B$. The SI factors are therefore not optional but
\emph{constitutive}: they are what makes the quantities involved physical
quantities in the first place. At the compact level there is neither $H_0$ nor
this bridge -- the question of the SI factor does not even arise there.

\subsection*{Verdict of the consistency check}

The bridge formula is \emph{consistent} -- but as a \emph{decompactified,
SI-anchored} relation, not as a compactified $T^4$ identity. The SI conversion
factor appears exactly as expected: partly cancelled ($\hbar c$, legitimate as
a length ratio), partly as the surviving empirical anchor
($m_e c^2/k_B T_\text{CMB}$), partly -- in the exponent notation -- as the
$11/2$ between the $E_0$ and Planck references \eqref{eq:exp-split}. This is
fully in line with Doc.~190: R34 (SI conversions indispensable for the
$c$-/time-power, $c=\hbar=1$ destroys the structure), R16 ($H_0$
emergent/decompactified) and R35 (no absorption of $\ell_P$ into the exponent).

% ============================================================
\section{Inversion: $H_0$ directly from the Casimir scale}
% ============================================================

Since $L_\xi$ follows via \eqref{eq:Lxi} from $\rho_\text{CMB}$ (i.e.\ from
$T_\text{CMB}$), \eqref{eq:bridge} can be solved for $H_0$:
\begin{equation}
  \boxed{\;H_0 = \frac{\pi}{2}\left[\frac{15}{\pi^2}\right]^{1/4}
  \frac{c}{L_\xi}\;\xi^{41/4}\;}
  \label{eq:h0fromcasimir}
\end{equation}
This is the experimental lever (Doc.~166, §5): a \emph{Casimir precision
measurement} of the scale $L_\xi$ fixes, via \eqref{eq:h0fromcasimir}, a value
for $H_0$ -- an $H_0$ determination in the laboratory rather than in the sky.

% ============================================================
\section{Complete number chain}
% ============================================================

\begin{center}
{\small
\renewcommand{\arraystretch}{1.3}
\adjustbox{max width=\textwidth}{%
\begin{tabular}{p{4.6cm}p{4.6cm}p{4.0cm}}
\toprule
\textbf{Quantity} & \textbf{Expression} & \textbf{Value} \\
\midrule
$\xi$                    & $4/30000$                              & $1{,}3333\times10^{-4}$ \\
$T_\text{CMB}$           & Planck 2018                            & $2{,}72548$\,K \\
$L_\xi$                  & $[15\xi/\pi^2]^{1/4}\,\hbar c/(k_BT_\text{CMB})$ & $100{,}24\,\mu$m \\
$\rho_\text{CMB}$        & $\xi\hbar c/L_\xi^4$                    & $4{,}170\times10^{-14}$\,J/m$^3$ \\
$\rho_\text{Casimir}/\rho_\text{CMB}$ & $\pi^2/(240\,\xi)$        & $\approx 308$ \\
$K_\text{geo}$           & $(\pi/2)(15/\pi^2)^{1/4}$              & $1{,}7441$ \\
$m_ec^2/(k_BT_\text{CMB})$ & --                                   & $2{,}176\times10^{9}$ \\
$\xi^{41/4}$             & $41/4 = 10{,}25$                       & $1{,}91\times10^{-40}$ \\
$\hbar H_0/(m_ec^2)$     & $(\pi/2)\,\xi^{10}$                    & $1{,}896\times10^{-38}$ \\
$L_\xi\cdot H_0/c$       & $K_\text{geo}\,(m_ec^2/k_BT_\text{CMB})\,\xi^{41/4}$ & $7{,}241\times10^{-31}$ \\
$H_0$                    & \eqref{eq:h0fromcasimir}               & $66{,}82$\,km/s/Mpc \\
\bottomrule
\end{tabular}%
}
}
\end{center}

% ============================================================
\section{Honest assessment of the status (Doc.~190)}
% ============================================================

The chain above is algebraically clean, but its \emph{epistemic} status must
remain clear (corrections register Doc.~190, entries R20, R32, R33, R35, R36,
R39):

\begin{warningboxenv}[What is derived -- and what is not]
\begin{itemize}
  \item \textbf{Form forward:} that the connection is a \emph{dimensionless
        $\xi$ ratio} (Casimir scale $\times$ Hubble scale $=$ $\xi$ power) is
        structurally forward. The \emph{chaining} of the building blocks
        \eqref{eq:bridgederiv} is exact.
  \item \textbf{Exponent not forward (R20):} the value $41/4$ -- already the
        partial exponent $10$ from \eqref{eq:h0} -- is \emph{not} derived
        forward from the T$^4$ geometry. The neat decomposition $41/4=10+1/4$
        \eqref{eq:decomp} is \emph{structurally suggestive}, not a proof.
  \item \textbf{$\xi^N$ guardrail (R35):} ``$X=\xi^N$'' counts as content only
        if $N$ is derived forward; otherwise any magnitude can be written as a
        $\xi$ power. The SI factor $\ell_P$ must not be tacitly absorbed into
        the exponent (cf.\ Place~3 above).
  \item \textbf{$H_0$ is emergent, not fundamental (R16/R32):} $H_0$ contains
        the $\hbar$ that becomes real only upon decompactification and does not
        exist in the compact $T^4$ form.
  \item \textbf{Shared reference point (R39):} the Planck length$\leftrightarrow
        R_H$ relation follows from \emph{no} framework out of first principles;
        $\Lambda$CDM fixes it via the measured $H_0$. FFGFT adopts the
        established value for comparability -- this is \emph{not} an asymmetric
        disadvantage but a shared empirical reference point of both models.
\end{itemize}
\end{warningboxenv}

\begin{foundation}[Verdict of the assessment]
The \emph{chaining} is exact -- a consistency relation among SI-anchored
quantities ($L_\xi$, $R_H$, $m_e$, $T_\text{CMB}$), together with the
experimental lever Casimir~$\to H_0$ \eqref{eq:h0fromcasimir}. $\ell_P$ does
\emph{not} appear in the formula; via the $E_0$ reference it is
\emph{extrapolable} -- as a consistency check, not a derivation: the leg
$E_0\leftrightarrow\ell_P$ ($11/2$) is derived (the fine-structure exponent,
see Outlook), the cosmic leg ($41/4$) remains open (R20). It is therefore
\emph{not} a first-principles prediction of $H_0$.
\end{foundation}

% ============================================================
\section{Outlook: the $11/2$, the $\Lambda$ fine-tuning problem and the $H_0$ tension}
% ============================================================

\textbf{The $11/2$ is a starting point, not a proof.} The exponent $11/2$ in
the split $63/4 = 41/4 + 11/2$ is not arbitrary: it coincides with the
\emph{forward-derived} fine-structure exponent,
\begin{equation}
  \alpha = \xi\cdot E_0^2 = K\cdot\xi^{11/2},\qquad E_0=\sqrt{m_e m_\mu}\propto\xi^{9/4}
  \quad\text{(Doc.~011, 070, 202, 210)},
\end{equation}
and it is the \emph{same} $E_0$ to which the cosmic exponent is referenced
($E_H=E_0\,\xi^{41/4}$, Doc.~026). Numerically $E_0$ sits exactly where the
split requires: $E_0/E_\text{Planck}=\xi^{5{,}48}$ versus $11/2=5{,}5$ --
a match to $\sim0{,}5\,\%$ (precisely the rounding $15{,}72$ vs.\ $15{,}75$
from R35).

\textbf{However:} the $\alpha$ exponent is referenced to the MeV / lepton base
(the $9/4$ scaling), the split $11/2$ to the Planck scale (the offset) --
different reference bases, the same number. Whether this is \emph{structurally
identical} or a numerical hit in the dense $\xi^N$ field is open; the guardrail
(R35, Doc.~271) requires for \enquote{content} an isolated, predicted, precise
value. The identity would be proven only by $E_0 = E_\text{Planck}\cdot\xi^{11/2}$
forward. \emph{Recorded as a conjecture, not as evidence.}

\textbf{The reverse question: why does $\Lambda$CDM's $H_0$ sit there?} Were
there a genuine, fixed connection $\ell_P\leftrightarrow R_H$, that would be a
sharp statement about $H_0$. But $\Lambda$CDM does \emph{not} choose $H_0$ from
theory -- it \emph{measures} it (and contested: the $H_0$ tension $\sim$67
vs.\ $\sim$73); $\ell_P$ comes from the measured $G,\hbar,c$. The ratio is not
harmless:
\begin{equation}
  \frac{\ell_P}{R_H}\approx 1{,}2\times10^{-61},
  \qquad
  \left(\frac{\ell_P}{R_H}\right)^2\approx10^{-122}
  \;\sim\;\frac{\rho_\Lambda}{\rho_\text{Planck}}.
\end{equation}
The $(\ell_P/R_H)^2\sim10^{-122}$ \emph{is} the cosmological constant problem --
the largest unexplained fine-tuning in physics. $\Lambda$CDM \enquote{chooses}
the value by fitting $\Lambda$ to the observed acceleration, without a
theoretical reason. FFGFT eliminates $\Lambda$ (R17/R39) and does \emph{not}
inherit the fine-tuning -- but relocates the same \enquote{why this scale} into
the exponent $41/4$.

\textbf{The test.} If the $41/4$ became forward-derivable (together with the
already-derived $11/2$ leg), FFGFT would \emph{predict} $\ell_P/R_H$ -- and
hence $H_0$ -- instead of importing it; it could even decide the $H_0$ tension
(which side, 67 or 73). Precisely that would turn the conjecture into evidence
(isolation/prediction/precision, Doc.~271). Until then it remains a suggestive
restatement of a measured ratio, open.
